\documentclass[11pt,a4paper]{article}

% --- Tectonic / XeTeX-safe preamble (no fontspec, no pandoc, no shell-escape) ---
\usepackage[margin=1in]{geometry}
\usepackage{lmodern}
\usepackage[T1]{fontenc}
\usepackage{microtype}
\usepackage{amsmath,amssymb}
\usepackage{graphicx}
\usepackage{booktabs}
\usepackage{array}
\usepackage{ragged2e}
\usepackage{enumitem}
\usepackage{xcolor}
\usepackage{titlesec}
\usepackage[font=small,labelfont=bf]{caption}
\usepackage[colorlinks=true,linkcolor=teal!70!black,citecolor=teal!70!black,urlcolor=teal!70!black]{hyperref}

\definecolor{ink}{RGB}{21,24,29}
\definecolor{accent}{RGB}{11,110,117}
\definecolor{good}{RGB}{46,115,80}
\definecolor{bad}{RGB}{156,51,46}
\definecolor{warn}{RGB}{143,98,16}

\titleformat{\section}{\Large\bfseries\color{ink}}{\thesection}{0.6em}{}
\titleformat{\subsection}{\large\bfseries\color{ink}}{\thesubsection}{0.6em}{}
\setlength{\parskip}{0.55em}
\setlength{\parindent}{0pt}
\newcommand{\code}[1]{\texttt{\small #1}}

\newcommand{\err}{\textcolor{bad}{\textbf{[error]}}}
\newcommand{\frm}{\textcolor{warn}{\textbf{[framing]}}}
\newcommand{\mth}{\textcolor{accent}{\textbf{[method]}}}

\title{\textbf{Open-Source Continuous-Culture Bioreactors:\\
A Correction, a Comparison, and a Two-Instrument Proposal}}
\author{Morgan Hough\\[2pt]
  \small Prepared with AI assistance (Claude, Anthropic).\\
  \small Source repository: \url{https://github.com/m9h/biopunk-bioreactor}}
\date{10 July 2026 \quad\textbullet\quad \small Working paper --- a reader's audit, not peer review}

\begin{document}
\maketitle

\begin{abstract}
\noindent
Low-cost open-source continuous-culture devices --- turbidostats, chemostats, and
morbidostats --- have proliferated, but they are frequently compared as though they compete
for a single role. They do not. We argue that the field serves \emph{three} distinct
purposes: in-situ \emph{characterisation} of a culture held in a defined state; automated
\emph{evolution} that must produce archived, replicated evidence; and turnkey
\emph{scale-out} operation. We first correct a set of verifiable errors and overstatements in
the recent OpenEvo preprint \cite{cocioba2026openevo}, including a self-contradictory
growth-rate figure and a false claim about a competitor's licensing. We then place six
platforms --- Chi.Bio, OpenEvo, Pioreactor, eVOLVER, EVE, and the Toprak morbidostat --- in a
like-for-like matrix across fifteen attributes, and give per-platform strengths and
limitations. Finally we propose that the best system is not one merged machine but
\emph{two} purpose-built instruments sharing a component ecosystem, and we note the strong
educational role these devices already play. Every factual claim is tagged with its
confidence level; specifications not verified against a primary source are flagged as such.
\end{abstract}

\section{Introduction}

Few organisms have shaped biology as thoroughly as brewer's yeast. The controlled fermentation
of \emph{Saccharomyces cerevisiae} is among the oldest of human technologies --- beer, bread,
and wine predate writing --- and for most of that history it was practised as a continuous or
serially propagated production process: a culture kept alive and productive across generations.
When that ancient craft finally met experimental science it did not merely benefit from biology;
it repeatedly \emph{founded} it. Pasteur's demonstration in the 1850s that fermentation is the
work of living yeast rather than spontaneous chemistry established microbiology and dealt the
decisive blow to spontaneous generation \cite{pasteur1857}. Four decades later, Buchner's finding
that a cell-free yeast extract still ferments sugar --- that the transformation is catalysed by
molecules, not a vital force --- founded biochemistry. The glycolytic pathway, the spine of
central metabolism, was mapped largely through fermentation; the first eukaryotic genome ever
sequenced was that of \emph{S.~cerevisiae} \cite{goffeau1996}; and the genetics of the cell
cycle, of protein secretion, and of homologous recombination were worked out in the same yeast
that raises bread.

Two threads of that history bear directly on this paper. The first is \emph{continuous culture}:
the mid-twentieth-century formalisation of steady-state growth --- Monod's growth kinetics
\cite{monod1949} and the chemostat of Novick and Szilard \cite{novick1950} --- turned the
fermentation vat into a scientific instrument,
a way to hold a population in a defined, unchanging environment and, crucially, to watch it mutate
and adapt under sustained selection. The second is \emph{production}: yeast
remains the workhorse eukaryotic cell factory, from industrial ethanol to the first recombinant
vaccine, and engineering better production strains is now a central problem of biotechnology.
Continuous culture and directed evolution are where these two threads meet --- an evolving
population, held in continuous culture, pushed toward a phenotype it did not begin with. The
open-source bioreactors reviewed here are the latest and cheapest instruments for exactly that
meeting, democratising a program that once demanded Pasteur's laboratory or an industrial
fermentation hall. It is fitting that OpenEvo's own optogenetic demonstration, and much of the
continuous-evolution work discussed below, is again carried out in yeast.

The precision of an in-vivo biology experiment depends on isolating the subsystem of interest
from the drift of a growing culture's environment. Continuous-culture devices address this by
diluting a culture during growth: a turbidostat holds optical density (OD) at a set-point; a
chemostat fixes the dilution rate and lets density float; a morbidostat titrates a drug to
hold growth inhibition constant. Over the past decade a wave of open-source, low-cost
instruments has made these techniques broadly accessible, and two in particular frame the
present analysis: \textbf{Chi.Bio} \cite{steel2020chibio}, an all-in-one platform for in-situ
measurement and manipulation, and \textbf{OpenEvo} \cite{cocioba2026openevo}, a recent
preprint describing a cheap, buildable device for automated directed evolution.

These two are often read as competitors --- both cost roughly \$300, both are open --- but
they are optimised for opposite halves of the problem. Chi.Bio is fundamentally an
\emph{instrument}: it exists to measure a culture richly while holding it in a defined state.
OpenEvo is fundamentally an \emph{actuator}: it exists to push a culture through selection
regimes, unattended, for weeks. Recognising this distinction --- and a third purpose that a
productised platform such as Pioreactor serves --- reorganises the entire field and changes
what ``best'' means.

This working paper has four aims. Section~\ref{sec:purposes} sets out the three purposes and
Section~\ref{sec:lineages} the founding literature and a landmark demonstration of each.
Section~\ref{sec:corrections} corrects specific errors in the OpenEvo preprint, fairly and
with proposed fixes. Section~\ref{sec:matrix} compares six platforms attribute by attribute,
and Section~\ref{sec:proposal} argues for building two purpose-built instruments rather than
one compromise machine. Section~\ref{sec:education} documents the educational materials these
platforms already serve. We are explicit throughout about what was verified against a primary
source and what was not.

\section{Three purposes, not one}
\label{sec:purposes}

Treating all continuous-culture devices as interchangeable obscures that they optimise for
incompatible objectives.

\textbf{\texttt{proto} --- Characterisation.} Hold a culture in a defined, unchanging state and
measure it richly, in situ, in real time, closing feedback loops around it. The figure of
merit is \emph{signal density per reactor}, with a handful of reactors running different
conditions in parallel. Chi.Bio is the exemplar: a chip spectrometer, a seven-colour LED for
optogenetic actuation and fluorophore excitation, non-contact infrared thermometry, and
proportional-integral-derivative plus model-predictive control that holds OD within about 2\%
of set-point \cite{steel2020chibio}.

\textbf{\texttt{devo} --- Evolution.} Push a culture through selection regimes unattended for
weeks, and produce \emph{evidence} that adaptation occurred and which mutation caused it. The
figures of merit are cheap chambers in numbers, systematic archiving, and long-run
robustness. This is OpenEvo's target, and --- as Section~\ref{sec:corrections} argues --- the
half it does not yet fully reach.

\textbf{\texttt{bior} --- Scale-out operation.} Run many reliable reactors now, with minimal build
effort, a maintained software platform, cluster management, and support. The value is
reproducibility and uptime, not measurement richness or archived evidence. Pioreactor
(turnkey) and eVOLVER (maximally configurable) own this axis, which is orthogonal to the other
two.

Read in order, the three trace a development lifecycle --- \texttt{proto} (prototype and
characterise a strain), \texttt{devo} (evolve it), \texttt{bior} (produce at scale) --- which
is why the tags are named for stages rather than numbered.

These purposes pull hardware in opposite directions --- rich-and-few versus cheap-and-many,
non-contact versus immersed-electrode-tolerant, defined-steady-state versus
free-to-run-to-stationary. A device optimised for one is measurably worse at another. That
tension is the organising principle of everything below.

\section{Three lineages: the founding literature of each purpose}
\label{sec:lineages}

The three tags are not conveniences of this review. Each names a distinct research tradition with
its own founding literature and its own landmark demonstration, and the traditions grew up
separately over some seventy years. Setting them side by side makes the central point concrete:
a single benchtop device now spans programs that were once the province of different laboratories
and different instruments.

\subsection{\texttt{proto} --- characterisation}
The characterisation tradition begins with the chemostat as a \emph{measurement} instrument.
Monod's growth kinetics and the Novick--Szilard chemostat \cite{monod1949,novick1950} made it
possible to hold a population at a defined, unchanging physiological state and read its properties
without the confounds of a changing batch environment. The decisive modern extension was to close
the loop --- to let a computer measure a living culture and act on it in real time; its founding
demonstration is Milias-Argeitis and colleagues' \emph{in silico} feedback, in which a controller
regulated a gene-expression circuit in vivo by continuously measuring output and adjusting an
input \cite{milias2011}. Chi.Bio \cite{steel2020chibio} later folded sensing, actuation, and
control into a single low-cost reactor.
\emph{Landmark demonstration:} automated optogenetic feedback control that held both gene
expression and growth rate at prescribed set-points in continuous \emph{E.~coli} culture, tracking
dynamic reference profiles and rejecting nutrient and temperature perturbations --- a bioreactor
regulating a cell's internal state the way a controller regulates an industrial plant
\cite{milias2016}.

\subsection{\texttt{devo} --- evolution}
The evolution tradition begins where continuous culture was first turned on selection: Novick and
Szilard used their chemostat not only to hold cells at steady state but to watch spontaneous
mutations sweep a bacterial population \cite{novick1950b}. It was systematised into experimental
evolution by Lenski's Long-Term Evolution Experiment, whose founding paper followed twelve
\emph{E.~coli} populations through their first 2{,}000 generations and established the discipline's
central apparatus --- replicate lineages and a frozen archive \cite{lenski1991}. For molecules
rather than whole organisms, the parallel founding is phage-assisted continuous evolution
\cite{esvelt2011pace}.
\emph{Landmark demonstration:} after roughly 31{,}500 generations one LTEE population evolved the
aerobic use of citrate --- a trait so unusual it helps define the species --- a result that could
be established, and its historical contingency proven by ``replaying'' the tape from frozen
ancestors, only because the lineages had been replicated and archived \cite{blount2012cit}.

\subsection{\texttt{bior} --- scale-out}
The scale-out tradition begins with continuous culture recast as a \emph{production} technology.
Herbert, Elsworth, and Telling's 1956 theoretical-and-experimental study established the
quantitative basis for running continuous culture as an industrial process and described an
operating pilot plant \cite{herbert1956}. Its modern form is parallelisation --- many small
reactors run at once --- from multiplexed chemostat arrays for high-throughput strain
characterisation \cite{miller2013} to the open, automated sixteen-vial eVOLVER
\cite{wong2018evolver}.
\emph{Landmark demonstration:} eVOLVER-hosted phage-assisted continuous evolution, in which
running many lagoons in parallel --- the throughput that scale-out exists to provide --- enabled
the evolution of compact Cas9 variants with new PAM specificities \cite{huang2022epace}.

These three lineages --- physiological measurement, experimental evolution, and production ---
matured in different laboratories with different apparatus. What is new is not any one of them but
that a device costing a few hundred dollars can now, in principle, serve all three at once; the
argument of this paper is precisely that it should not try to.

\section{Corrections to the OpenEvo preprint}
\label{sec:corrections}

OpenEvo is a genuine contribution and the corrections here are not a dismissal. Its real
strengths include a peak-to-peak growth readout tied to hardware dilution events (which
sidesteps local-maxima noise); a 940\,nm OD path that avoids optogenetic and ambient overlap,
with an inverse model holding $R^2 \ge 0.98$ out to $\mathrm{OD}_{600}=6.0$; a closed
autoclavable vessel with better sterility than Chi.Bio's open tube; a 4{,}000-hour remote run
demonstrating robustness; and an illustrated assembly manual validated by having an
inexperienced undergraduate build a working unit. It is the cheapest buildable three-media
selection cycler in the field. The issues below are correctable and concern framing and
analysis more than hardware.

\subsection{Verifiable errors}

\err\ \textbf{The headline growth-rate figure contradicts itself.} The abstract states
evolved cells ``grew 55\% faster than wild-type at 12\% salt''; the results state ``a doubling
time of 4.9 hours versus 7.6 hours (a 36\% increase in growth rate).'' Both describe the same
data, but the two numbers are different quantities. A doubling time of $7.6 \to 4.9$\,h is a
36\% \emph{reduction in doubling time} yet a 55\% \emph{increase in growth rate}, since rate
$\propto 1/(\text{doubling time})$ and $(1/4.9)/(1/7.6) = 1.55$. The abstract's 55\% is the
correct growth-rate figure; the results sentence mislabels the 36\% doubling-time reduction as
a growth-rate increase. \emph{Fix:} state both explicitly --- ``36\% shorter doubling time,
equivalently 55\% higher growth rate.''

\err\ \textbf{A false claim about a competitor.} The preprint states that Pioreactor is
``available for purchase, but only the software is open source.'' A public
\code{Pioreactor/hardware} repository exists and the hardware is openly licensed under
CC~BY-SA~4.0 (per its \code{LICENSE} file) \cite{pioreactor}. The claim is used to clear
competitive space and is false. \emph{Fix:} describe Pioreactor accurately as
open-hardware-and-software, distinguished not by openness but by being a supported product
with a fleet platform. (We note in passing that an earlier draft of this document itself
misnamed the licence as CERN-OHL; a contributor corrected it against the primary
\code{LICENSE} file --- an illustration of the citation discipline this repository aims for.)

\subsection{Framing and omission}

\frm\ \textbf{The nearest prior art is uncited.} Chi.Bio \cite{steel2020chibio} is open
hardware and software, roughly the same \$300, performs in-situ measurement and continuous
culture, and predates OpenEvo by six years. The preprint surveys eVOLVER, the Toprak
morbidostat, and Pioreactor, then frames an unfilled gap while omitting the system that most
directly occupies it, which makes the gap look larger than it is.

\frm\ \textbf{An unsupported superlative.} ``The most affordable modern continuous growth
system available'' is not established: Pioreactor ships turnkey at \$329 with no build, and EVE
builds a triplicate system for \$115--200 \cite{gopalakrishnan2022eve}. \emph{Fix:} qualify to
``among the lowest-cost \emph{buildable} systems with three-media selection.''

\frm\ \textbf{The success is overstated.} The 9\% salt step did not partially adapt --- growth
ceased, OD flat-lined, and the run terminated at cycle 152. The adaptation claim rests on a
single hand-picked glycerol stock at cycle 145, assayed offline. The result is real, but the
framing should foreground one successful selection step with a failed extension, not a ladder.

\subsection{Methodological limitations}

\mth\ \textbf{$n=1$ forecloses causal attribution.} A single clone from a single population was
sequenced. Without independent replicate lineages there is no convergence signal, and the
reported single-nucleotide polymorphisms and two large deletions cannot be attributed to the
low-salt adaptation. This is the field's central inference requirement --- the Long-Term
Evolution Experiment runs twelve populations for exactly this reason
\cite{blount2012cit} --- and it is the most consequential limitation.

\mth\ \textbf{The wrong variant caller.} Structural variants were called with ``Geneious's
default mapper with low-medium sensitivity for 5 iterations.'' The key finding --- a
$\sim$110\,kb pHV4 deletion with ``primarily transposon-associated elements'' --- is
IS-mediated rearrangement, precisely the case a default short-read mapper handles worst.
\code{breseq} exists for microbial-evolution resequencing including IS mobilisation
\cite{deatherage2014breseq}, and its author is cited five times in the preprint. \emph{Fix:}
re-call with \code{breseq} and confirm the pHV4 structure with long reads.

\mth\ \textbf{No systematic archiving.} The preprint cites Barrick et al.\ ---
``\dots Archiving Populations, and Measuring Fitness\dots'' \cite{barrick2023ltee} --- as
motivation, then implements neither archiving nor fitness measurement in the automated system.
Head-to-head ancestor-versus-descendant competition, the gold standard the cited work
describes, is impossible without a time-series archive.

\section{A like-for-like comparison}
\label{sec:matrix}

Table~\ref{tab:matrix} places six platforms across fifteen attributes. Purpose codes are
proto~(characterise), devo~(evolve), bior~(scale-out). Cells reflect each platform as published or
shipped, not what a motivated builder could add; ``independent use'' means peer-reviewed work
by groups other than the originators.

\begin{table}[t]
\centering
\footnotesize
\renewcommand{\arraystretch}{1.25}
\caption{Six open continuous-culture platforms, attribute by attribute.
\textcolor{good}{Green} marks a notable strength, \textcolor{bad}{red} a notable weakness,
\textcolor{warn}{amber} a qualified case.}
\label{tab:matrix}
\begin{tabular}{@{}>{\RaggedRight}p{2.4cm} p{1.85cm} p{1.85cm} p{1.85cm} p{1.85cm} p{1.65cm}@{}}
\toprule
\textbf{Attribute} & \textbf{Chi.Bio} & \textbf{OpenEvo} & \textbf{Pioreactor} & \textbf{eVOLVER} & \textbf{EVE} \\
\midrule
Primary purpose & proto & devo & bior & bior & devo \\
Reference & PLoS Biol '20 & bioRxiv '26 & \textcolor{bad}{none} & Nat Biot '18 & eLife '22 \\
Parallelism & \textcolor{good}{8 / computer} & \textcolor{bad}{1 chamber} & \textcolor{good}{cluster} & \textcolor{good}{16 vials} & $\ge$3 reps \\
Working vol. & 12--25 mL & $\sim$20 mL & 15--30 mL & $\sim$40 mL & $\sim$12 mL \\
OD source & 650 nm laser & 940 nm IR & 900 nm IR & LED/PD & LED/PD \\
Rich sensing & \textcolor{good}{spectrometer} & \textcolor{bad}{OD only} & plugin & config. & \textcolor{bad}{OD only} \\
Contact & \textcolor{good}{non-contact} & closed & contact & contact & contact \\
Sterility & \textcolor{warn}{open air} & \textcolor{good}{closed, autocl.} & autocl. vial & autocl. vial & 3D-print \\
Selection & 2+2 pumps & \textcolor{good}{3 media+waste} & pumps & \textcolor{good}{config.} & drug line \\
Control law & \textcolor{good}{PID + MPC} & threshold & Kalman GR & config. & threshold \\
Archiving & \textcolor{bad}{none} & \textcolor{bad}{manual 1$\times$} & \textcolor{bad}{none} & \textcolor{bad}{none} & \textcolor{bad}{none} \\
Open HW/SW & \textcolor{good}{both} & \textcolor{good}{both} & \textcolor{good}{both} & \textcolor{good}{both} & \textcolor{good}{both} \\
Cost/unit & \$300+PCB & $\sim$\$300 & \$329--389 & \textcolor{warn}{$<$\$2k/16} & \textcolor{good}{\$115--200} \\
Build effort & \textcolor{warn}{PCB fab} & \textcolor{good}{guided} & \textcolor{good}{turnkey} & \textcolor{bad}{machine shop} & moderate \\
Independent use & ReacSight & \textcolor{warn}{too new} & \textcolor{warn}{community} & \textcolor{good}{ePACE} & teaching \\
\bottomrule
\end{tabular}
\end{table}

The Toprak morbidostat \cite{toprak2013morbidostat} is omitted from the table for width but
belongs to purpose devo: it is the canonical antibiotic-resistance device (working volume
$\sim$12\,mL), titrating a drug to constant inhibition, widely adapted but a dated bespoke
build rather than a maintained platform.

Two patterns are worth naming. First, \emph{every platform lacks archiving} --- the single
most important capability for turning an evolution run into evidence, and a clear gap for the
whole field. Second, the only platform without a reference paper is Pioreactor, which is also
the only pure bior product; its credibility is commercial and community rather than
bibliographic, a genuine consideration when an instrument must be cited.

\section{Why the evolution slot has no champion}
\label{sec:devo}

Two of the three purposes have clear champions --- Chi.Bio for \texttt{proto}, Pioreactor and
eVOLVER for \texttt{bior} --- but the \texttt{devo} slot is conspicuously empty, and OpenEvo,
the one device built expressly for it, does not fill it. The reason is definitional.
\texttt{devo} is not ``push a culture through a selection regime'' but ``produce an archived,
replicated, sequenced record that establishes that adaptation occurred and identifies the
causal mutation.'' A champion must deliver the purpose's output, and OpenEvo delivers the run,
not the evidence.

The shortfall is exactly the set of limitations catalogued in Section~\ref{sec:corrections}.
With a single chamber and $n=1$ there are no replicate lineages and hence no convergence signal,
so the observed mutations cannot be attributed to selection rather than to drift or hitchhiking
\cite{blount2012cit}. With no systematic archiving there is no frozen fossil record, and without
one the ancestor-versus-descendant fitness competition that makes evolution quantitative is
impossible \cite{barrick2023ltee}. And the single structural variant reported was called with a
default short-read mapper rather than \code{breseq} \cite{deatherage2014breseq}. OpenEvo is, in
short, an excellent actuator with no evidentiary apparatus attached --- the best pipette in the
room with no freezer beside it.

Two further considerations keep the title vacant. The first is adoption: a champion earns
standing through independent use by groups other than its originators, and the incumbent
champions have it. Chi.Bio has been taken up by \emph{ReacSight}, a framework that couples open
bioreactor arrays --- Chi.Bio among them --- to an automated flow cytometer and a liquid-handling
robot, so that experiments run \emph{reactively}: each measurement feeds back to change the next
actuation, turning a passive array into a closed-loop platform \cite{fribourg2022reacsight}.
eVOLVER, for its part, has become a substrate for continuous directed evolution: it has been
used to host \emph{phage-assisted continuous evolution} (PACE) --- the method of Esvelt,
Carlson, and Liu in which a desired molecular activity is linked to the propagation of an M13
bacteriophage, so that proteins evolve tens of generations per day with minimal human
intervention \cite{esvelt2011pace} --- run across many parallel lagoons rather than the single
bespoke apparatus of the original, a throughput increase that directly enabled campaigns such as
the evolution of compact Cas9 variants \cite{huang2022epace}. This automated, eVOLVER-hosted form
is what the community means by \emph{ePACE}. OpenEvo, days old at this writing, has no independent
use yet.

The second consideration is more pointed: the remedy does not obviously return to OpenEvo. What
would constitute a \texttt{devo} champion --- of order eight replicate chambers under one
coordinating interface, an automated chilled archiver, and a \code{breseq} pipeline --- most
closely resembles Pioreactor's existing leader/worker cluster with an archiving plug-in, not a
scaled-up OpenEvo. OpenEvo's durable advantage then narrows to its three-media selection fluidics
and its closed autoclavable vessel, both of which are graftable onto another chassis.

None of this diminishes what OpenEvo is: the cheapest buildable three-media selection cycler in
the field, with the best assembly manual among the six platforms, and the leading \emph{candidate}
for the role. But candidate is not champion, and the distance between them is precisely archiving,
replication, and correct variant calling. Close those three and the slot is filled.

\subsection{What the actuator must be connected to}
\label{sec:apparatus}

The corollary of defining \texttt{devo} by evidence is that a champion is not a better actuator
but an actuator wired into an evidentiary pipeline. OpenEvo already performs the
actuation --- controlled dilution, three-media selection, unattended runs of thousands of
hours. What it lacks is everything \emph{downstream} of the culture vessel. Five connections
turn the actuator into a champion, and only one of them is new hardware.

\textbf{1. A replicate manifold.} The actuator must be one of many identical chambers under a
single coordinating controller, so that a campaign runs of order eight or more independent
lineages in parallel rather than one. Convergence across those lineages --- the same locus
mutated repeatedly --- is the signal that a mutation was selected rather than drifting
\cite{blount2012cit}, and it cannot be reconstructed after the fact from a single run. This
connection is a fleet interface, and the reference implementation already exists as Pioreactor's
leader/worker cluster; the actuator connects \emph{to} it rather than reinventing it.

\textbf{2. A cold archive.} The vessel's outflow --- which already carries culture out on every
dilution --- must connect to an automated chilled fraction collector that banks a glycerol stock
at programmed generation intervals. This is the frozen fossil record, the single component whose
absence most distinguishes every current open platform from the Long-Term Evolution Experiment
\cite{barrick2023ltee}, and the one genuinely new module the champion requires. It is
inexpensive: a refrigerated carousel and a diverter valve on an existing pump line.

\textbf{3. A sequencing and variant-calling pipeline.} Archived samples must flow to
whole-genome sequencing and be called with \code{breseq} \cite{deatherage2014breseq}, whose
handling of insertion-sequence--mediated structural variation is exactly what a default mapper
lacks. This connection is off-instrument --- a sample-submission workflow and a compute
pipeline --- but it is part of the apparatus, because an evolved strain whose mutations are
miscalled is not evidence.

\textbf{4. A fitness-measurement loop.} The archive closes back onto the actuator: reviving a
banked ancestor and competing it head-to-head against a descendant, read out through the same
optical-density hardware, is what converts ``the population changed'' into a measured selection
coefficient. Without the archive there is nothing to revive; with it, the actuator's own sensors
become the fitness assay, and no new instrument is needed.

\textbf{5. A selection-design layer.} Finally, the actuator must be connected to a principled
choice of selection pressure. Naive growth selection on a production strain evolves the mutant
that deletes the costly pathway \cite{sandberg2019ale}; the champion couples product formation
to growth --- or deploys OpenEvo's negative-selection media line as counterselection --- before
the run begins. This connection is upstream and computational, a genome-scale design step, but it
determines whether the evidence the apparatus so carefully collects is evidence of anything worth
having.

The shape of the argument is therefore that the \texttt{devo} champion is mostly
\emph{connective tissue}: one new hardware module (the archiver), one borrowed fleet interface,
one off-instrument pipeline (sequencing plus \code{breseq}), one feedback loop that reuses
existing sensors, and one upstream design step. The actuator itself stays cheap and
simple --- OpenEvo's genuine strength --- and the championship is earned by what surrounds it,
not by loading the vessel with sensors. This is the same conclusion the two-instrument proposal
reaches from the opposite direction (Section~\ref{sec:proposal}): resist instrumenting the
chamber, and invest in the apparatus around it.

\section{A two-instrument proposal}
\label{sec:proposal}

Because the three purposes pull hardware in opposite directions, the strongest system is not
one merged machine but two purpose-built instruments that share a component ecosystem.

\subsection{The best characterisation instrument}
Start from Chi.Bio; its architecture is already correct. The governing constraint is that
\emph{non-contact measurement is inviolable} --- every subsystem but the heat plate reads
through the vessel wall, which is why tubes hot-swap and sterilisation is trivial. Upgrades
that respect this: dissolved-oxygen and pH optode spots read optically through the wall; a
dispersive micro-spectrometer replacing eight fixed filter bands with a continuous
340--850\,nm trace; and a sealed, gas-controlled lid to close Chi.Bio's open-to-atmosphere
gap. Viable-cell-density by dielectric spectroscopy is explicitly \emph{excluded} here because
it requires immersed electrodes --- it belongs on the other machine.

\subsection{The best evolution machine}
Start from OpenEvo; its virtues (cheap, closed, buildable, standalone) fit long unattended
runs. But its two missing capabilities are not sensors. First, \emph{replicate populations}:
$n=1$ is a failure of causal inference, so the build is eight cheap chambers under one
coordinating interface, not one richly instrumented reactor. Second, \emph{archiving}: an
automated chilled fraction collector on the existing waste line, diverting a programmable
fraction on a generation-count schedule into a chilled glycerol archive --- the capability
that makes fitness measurable, and, as far as we can determine, one nobody has built. A
\emph{batch-excursion} firmware mode --- periodically suspending dilution to run to stationary
phase --- would additionally recover carrying capacity, a phenotype a turbidostat structurally
cannot observe, and the very axis on which the OpenEvo cycle-145 population differed. The
adaptive-laboratory-evolution literature adds a crucial caution: selection is for growth, and
product formation costs growth, so growth-coupled strain design must precede selection
\cite{sandberg2019ale}.

\subsection{Why two, not one}
The sensor modalities sort themselves by contact. Optode DO/pH and dispersive spectrometry are
optical and non-contact, and land on the characterisation instrument; conductivity, electrical
resistance tomography, and dielectric viable-cell-density need immersed electrodes and land on
the evolution machine, whose closed autoclavable vessel welcomes them. Almost no sensor is
wanted by both and available to only one --- the clearest evidence these are two instruments.
What they \emph{share} is an ecosystem: a common module footprint, reusable sensor-module
designs, one firmware framework, and a shared calibration reference. Combine at the platform
level, not the product level.

\section{Educational role}
\label{sec:education}

These platforms are teaching instruments as much as research tools, and because one device
spans microbiology, electronics, firmware, control theory, and data analysis, a single unit
serves biology, engineering, and computing classes at once. EVE is the education-first
exemplar: its paper is explicitly titled ``\dots and its educational use,'' it builds for under
\$200, and it was validated in a real classroom evolving replicate \emph{E.~coli} under
chloramphenicol \cite{gopalakrishnan2022eve}. Pioreactor maintains an educators portal and
lesson concepts spanning OD-by-scattering through evolution. OpenEvo's assembly manual doubles
as a build-it-yourself course and anchored an international hackathon in which students
remotely controlled a device on another continent. Chi.Bio's own documentation offers its core
features as a stand-alone teaching platform on which students can observe fluorescence by eye.
Notably, the teaching materials inherit the research blind spots: archiving, replicate-based
inference, and carrying capacity are as under-taught as they are under-implemented --- an
opportunity for a device-agnostic curriculum.

\section{Discussion and limitations}
\label{sec:discussion}

This is a reader's audit, not peer review, and confidence is uneven. Full texts of Chi.Bio and
OpenEvo were read directly. Pioreactor specifications, pricing, hardware licensing
(CC~BY-SA~4.0, verified against the primary \code{LICENSE} file), and the absence of a
reference paper were confirmed from primary or product sources, as were the headline
specifications of eVOLVER and EVE and the working volume of the Toprak morbidostat. Component
datasheets underpinning the proposed evolution machine (spectrometer, spectral sensor, and
impedance front-ends) were verified in a companion pass. Some adoption specifics and a small
number of vendor register-level details remain unverified at page level and are flagged as
such in the repository's \code{SOURCES.md}. No claim here should be budgeted against without
confirmation at its primary source.

\section{Conclusion}
\label{sec:conclusion}

The open-bioreactor field is best understood as serving three purposes, not one contest.
Chi.Bio owns characterisation; Pioreactor and eVOLVER own scale-out; and evolution --- the
purpose with the most devices chasing it --- has no clear champion, because no platform yet
archives or replicates as the science requires. OpenEvo's real contribution is the cheapest
buildable selection cycler with a genuinely excellent manual; its preprint's errors of framing
and analysis are correctable and worth correcting. The best path forward is two purpose-built
instruments sharing a parts ecosystem, with archiving and replication treated as
first-class --- and with the educational mission, already strong, extended by a curriculum that
teaches the methodological rigour the hardware has so far left out.

\vspace{1em}
\hrule
\vspace{0.5em}
{\small\textit{Source data, the full comparison matrix, per-platform pros and cons, a
datasheet-verification pass, an educational-materials guide, and this document's \LaTeX{}
sources are maintained at} \url{https://github.com/m9h/biopunk-bioreactor} \textit{under
CC~BY~4.0. Corrections with citations are welcome.}}

\begin{thebibliography}{10}\small
\bibitem{steel2020chibio} Steel H, Habgood R, Kelly CL, Papachristodoulou A.
  In situ characterisation and manipulation of biological systems with Chi.Bio.
  \textit{PLoS Biology} 18(7):e3000794 (2020).
\bibitem{cocioba2026openevo} Cocioba SS, Huang P-C, Mallon J, Chan Z, Geremew AW, Bisson A,
  Kyriakakis P. OpenEvo: An Open-Source Platform for Automated Evolution and Analysis.
  \textit{bioRxiv} 2026.07.06.735356 (2026). Preprint.
\bibitem{wong2018evolver} Wong BG, Mancuso CP, Kiriakov S, Bashor CJ, Khalil AS.
  Precise, automated control of conditions for high-throughput growth of yeast and bacteria
  with eVOLVER. \textit{Nature Biotechnology} 36(7):614--623 (2018).
\bibitem{gopalakrishnan2022eve} Gopalakrishnan V, et al., Scott JG.
  A low-cost, open-source evolutionary bioreactor and its educational use.
  \textit{eLife} 11:e83067 (2022).
\bibitem{toprak2013morbidostat} Toprak E, Veres A, Yildiz S, Pedraza JM, Chait R, Paulsson J,
  Kishony R. Building a morbidostat. \textit{Nature Protocols} 8(3):555--567 (2013).
\bibitem{pioreactor} Davidson-Pilon C. Pioreactor: an accessible, extensible bioreactor
  platform. \url{https://pioreactor.com} (2023). CC~BY-SA~4.0 hardware, MIT software; no
  reference paper.
\bibitem{barrick2023ltee} Barrick JE, et al. Daily Transfers, Archiving Populations, and
  Measuring Fitness in the LTEE with \textit{E. coli}. \textit{J Vis Exp} (2023).
\bibitem{blount2012cit} Blount ZD, Barrick JE, Davidson CJ, Lenski RE. Genomic analysis of a
  key innovation in an experimental \textit{E. coli} population. \textit{Nature} 489:513--518 (2012).
\bibitem{sandberg2019ale} Sandberg TE, Salazar MJ, Weng LL, Palsson BO, Feist AM. The emergence
  of adaptive laboratory evolution\dots \textit{Metabolic Engineering} 56:1--16 (2019).
\bibitem{deatherage2014breseq} Deatherage DE, Barrick JE. Identification of mutations in
  laboratory-evolved microbes\dots using breseq. \textit{Methods Mol Biol} 1151:165--188 (2014).
\bibitem{fribourg2022reacsight} Bertaux F, Aditya C, et al. Enhancing bioreactor arrays\dots
  with ReacSight. \textit{Nature Communications} 13:3363 (2022).
\bibitem{esvelt2011pace} Esvelt KM, Carlson JC, Liu DR. A system for the continuous directed
  evolution of biomolecules. \textit{Nature} 472:499--503 (2011).
\bibitem{huang2022epace} Huang TP, Heins ZJ, et al. High-throughput continuous evolution of
  compact Cas9 variants (eVOLVER-hosted PACE). \textit{Nature Biotechnology} 40:1918--1930 (2022).
\bibitem{pasteur1857} Pasteur L. M\'emoire sur la fermentation alcoolique.
  \textit{Comptes Rendus de l'Acad\'emie des Sciences} 45:1032--1036 (1857).
\bibitem{novick1950} Novick A, Szilard L. Description of the chemostat.
  \textit{Science} 112:715--716 (1950).
\bibitem{monod1949} Monod J. The growth of bacterial cultures.
  \textit{Annual Review of Microbiology} 3:371--394 (1949).
\bibitem{goffeau1996} Goffeau A, Barrell BG, Bussey H, et al. Life with 6000 genes.
  \textit{Science} 274:546--567 (1996).
\bibitem{milias2011} Milias-Argeitis A, Summers S, Stewart-Ornstein J, Zuleta I, Pincus D,
  El-Samad H, Khammash M, Lygeros J. In silico feedback for in vivo regulation of a gene
  expression circuit. \textit{Nature Biotechnology} 29:1114--1116 (2011).
\bibitem{milias2016} Milias-Argeitis A, Rullan M, Aoki SK, Buchmann P, Khammash M. Automated
  optogenetic feedback control for precise and robust regulation of gene expression and cell
  growth. \textit{Nature Communications} 7:12546 (2016).
\bibitem{novick1950b} Novick A, Szilard L. Experiments with the Chemostat on spontaneous
  mutations of bacteria. \textit{PNAS} 36(12):708--719 (1950).
\bibitem{lenski1991} Lenski RE, Rose MR, Simpson SC, Tadler SC. Long-term experimental evolution
  in \textit{Escherichia coli}. I. Adaptation and divergence during 2,000 generations.
  \textit{The American Naturalist} 138(6):1315--1341 (1991).
\bibitem{herbert1956} Herbert D, Elsworth R, Telling RC. The continuous culture of bacteria: a
  theoretical and experimental study. \textit{Journal of General Microbiology} 14(3):601--622 (1956).
\bibitem{miller2013} Miller AW, Befort C, Kerr EO, Dunham MJ. Design and use of multiplexed
  chemostat arrays. \textit{Journal of Visualized Experiments} (72):e50262 (2013).
\end{thebibliography}

\end{document}
